Algorithmic redistricting depends on three data sources: the Census
P.L.\ 94-171 redistricting file, TIGER/Line shapefiles, and optionally
American Community Survey 5-year tables.
A hostile actor who can modify any of these sources before the algorithm
processes them can potentially steer partisan outcomes without touching
algorithm code or parameters.

We analyse the three attack surfaces---population count manipulation,
shapefile boundary modification, and ACS demographic table alteration---
and show that (1)~even a $\pm 1\%$ per-tract population manipulation
produces at most $\Delta D \leq 0.4$ expected Democratic seats nationally,
placing the effect within the B.7 seed-variance noise floor;
(2)~all three attack surfaces are protected by \bisect{}'s three-level
SHA-256 audit chain: Census download hash, adjacency construction hash,
and final plan hash; and (3)~any manipulation that exceeds the noise floor
requires modifications large enough to produce detectable hash mismatches.

A disclosed limitation: \bisect{} uses TLS certificate validation via
the system trust store but does not implement certificate pinning.
A sophisticated man-in-the-middle attacker with TLS spoofing capability
could intercept Census downloads without changing the algorithm's behavior.
We recommend supplementary out-of-band hash verification for litigation-grade deployments.

The three-level audit chain---implemented via \texttt{bisect fetch}
and \texttt{bisect label-verify}---provides a computationally infeasible
barrier to undetected manipulation: SHA-256 collision resistance
requires $2^{80}$ operations under the birthday bound
(FIPS PUB 180-4~\citep{fips180}).
